\section{Aritmetika Matriks}

\begin{outcome}
\begin{enumerate}
\item[A.] Melakukan operasi matriks berupa penjumlahan matriks, perkalian
skalar, transpos, dan perkalian matriks. Mengidentifikasi kapan
operasi-operasi ini tidak terdefinisi. Merepresentasikan operasi-operasi ini berdasarkan
entri-entri suatu matriks.

\item[B.] Membuktikan sifat-sifat aljabar untuk penjumlahan matriks, perkalian
skalar, transpos, dan perkalian matriks. Menerapkan sifat-sifat ini
untuk memanipulasi ekspresi aljabar yang melibatkan matriks.

\item[C.] Menghitung invers suatu matriks menggunakan operasi baris, dan membuktikan identitas yang melibatkan
invers matriks.

\item[E.] Menyelesaikan sistem linear menggunakan aljabar matriks.

\item[F.] Menggunakan perkalian dengan matriks elementer untuk menerapkan operasi baris. 

\item[G.] Menuliskan suatu matriks sebagai hasil kali matriks-matriks elementer.

\end{enumerate}
\end{outcome}

Anda kini telah menyelesaikan sistem persamaan dengan menuliskannya dalam bentuk
matriks diperbesar, lalu melakukan operasi baris pada matriks diperbesar tersebut. Ternyata,
matriks penting bukan hanya untuk sistem persamaan, melainkan juga dalam banyak penerapan.

Ingat bahwa sebuah \textbf{matriks}
\index{matriks} adalah susunan bilangan berbentuk persegi panjang. Bentuk jamaknya juga
disebut \textbf{matriks}. Sebagai contoh, berikut adalah sebuah matriks.
\begin{equation}
\leftB
\begin{array}{rrrr}
1 & 2 & 3 & 4 \\
5 & 2 & 8 & 7 \\
6 & -9 & 1 & 2
\end{array}
\rightB
\label{matrix}
\end{equation}
Ingat bahwa ukuran atau dimensi suatu matriks didefinisikan sebagai $m\times n$, dengan $m$ sebagai
banyaknya baris dan $n$ sebagai banyaknya kolom. Matriks di atas adalah matriks
$3\times 4$ karena memiliki tiga baris dan empat kolom. Anda dapat mengingat bahwa kolom-kolom
menyerupai tiang pada kuil Yunani. Kolom berdiri tegak, sedangkan baris membentang
mendatar seperti alur yang dibuat traktor di ladang yang telah dibajak.

Saat menyatakan ukuran matriks, banyaknya baris selalu ditulis
sebelum banyaknya kolom. Anda
dapat mengingat urutan ini dengan frasa \textbf{B}aris sebelum \textbf{K}olom. 

Perhatikan definisi berikut.

\begin{definition}{Matriks Persegi}\label{def:squarematrix}
Matriks $A$ yang berukuran $n \times n$ disebut \textbf{matriks persegi} \index{matriks!persegi}.
Dengan kata lain, $A$ adalah matriks persegi jika banyaknya baris sama dengan
banyaknya kolom.
\end{definition}

Sekarang kita memperkenalkan beberapa notasi yang khusus digunakan untuk matriks. Kita menyatakan kolom-kolom matriks $A$
dengan $A_{j}$ sebagai berikut
\begin{equation*}
A = 
\leftB
\begin{array}{rrrr}
A_{1} & A_{2} & \cdots & A_{n}
\end{array}
\rightB
\end{equation*}
Dengan demikian, $A_{j}$ adalah kolom $j$ dari $A$, jika dihitung dari kiri ke kanan. 

Unsur-unsur individual dalam matriks disebut \textbf{entri} \index{matriks!entri matriks} atau \textbf{komponen} \index{matriks!komponen matriks} dari $A$. Unsur-unsur matriks
diidentifikasi berdasarkan posisinya. $\mathbf{\left( i, j \right)}$\textbf{-entri} suatu matriks adalah entri
pada baris $i$ dan kolom $j$. Sebagai contoh, pada matriks \ref{matrix} di atas, $8$ berada di
posisi $\left(2,3 \right)$ (dan disebut entri-$\left(2,3 \right)$) karena berada pada baris kedua dan kolom ketiga. 

Agar jelas matriks mana yang sedang dibahas, kita
akan menyatakan entri pada baris $i$ dan kolom $j$ dari matriks $A$ dengan $a_{ij}$. Dengan demikian, kita dapat menuliskan $A$ berdasarkan entri-entrinya,
sebagai $A= \leftB a_{ij} \rightB$. Dengan menggunakan notasi ini pada matriks dalam \ref{matrix},
$a_{23}=8, a_{32}=-9, a_{12}=2,$ dan seterusnya.

Ada berbagai operasi yang dilakukan pada matriks-matriks dengan ukuran yang sesuai.
Matriks dapat dijumlahkan dengan dan dikurangkan dari matriks lain,
dikalikan dengan suatu skalar, dan dikalikan dengan matriks lain. Kita
tidak pernah membagi suatu matriks dengan matriks lain, tetapi nanti kita akan melihat bagaimana invers matriks memainkan peran yang serupa. 

Dalam melakukan aritmetika matriks, kita sering mendefinisikan suatu operasi berdasarkan apa yang
terjadi pada entri-entri (atau komponen-komponen)
matriks tersebut. Sebelum menelaah operasi-operasi ini secara mendalam, perhatikan beberapa
definisi umum berikut.

\begin{definition}{Matriks Nol}\label{def:zeromatrix}
\textbf{Matriks nol $m\times n$} adalah matriks $m\times n$
yang setiap entrinya sama dengan nol. Matriks ini
\index{matriks nol} dinyatakan dengan $0.$
\end{definition}

Salah satu matriks nol ditunjukkan dalam contoh berikut.

\begin{example}{Matriks Nol}\label{exa:zeromatrix}
Matriks nol $2\times 3$ adalah $0= \leftB
\begin{array}{ccc}
0 & 0 & 0 \\
0 & 0 & 0
\end{array}
\rightB $.
\end{example}

Perhatikan bahwa ada matriks nol $2\times 3$, matriks nol $3\times 4$, dan seterusnya. Bahkan,
ada matriks nol untuk setiap ukuran! 

\begin{definition}{Kesamaan Matriks}\label{def:equalityofmatrices}
 Misalkan $A$ dan $B$ adalah dua matriks $m \times n$. Maka $A=B$ \index{matriks!kesamaan} berarti
bahwa untuk $A=\leftB a_{ij}\rightB $
dan $B=\leftB b_{ij}\rightB ,$ berlaku $a_{ij}=b_{ij}$ untuk semua $1\leq i\leq m$ dan
$1\leq j\leq n.$
\end{definition}

Dengan kata lain, dua matriks sama persis ketika ukurannya sama dan
entri-entri yang bersesuaian identik. Oleh karena itu,
\begin{equation*}
\leftB
\begin{array}{rr}
0 & 0 \\
0 & 0 \\
0 & 0
\end{array}
\rightB \neq \leftB
\begin{array}{rr}
0 & 0 \\
0 & 0
\end{array}
\rightB
\end{equation*}
karena ukurannya berbeda. 
Selain itu,
\begin{equation*}
\leftB
\begin{array}{rr}
0 & 1 \\
3 & 2 
\end{array}
\rightB \neq \leftB
\begin{array}{rr}
1 & 0 \\
2 & 3
\end{array}
\rightB
\end{equation*}
karena, meskipun ukurannya sama, entri-entri yang bersesuaian tidak identik.

Pada bagian berikutnya, kita membahas penjumlahan matriks. 
